\chapter{Aritmetika}\label{ch:g12:arith}

Aritmetika mengkaji \omterm{def:g10:numbers:sets}{bilangan bulat}: \omterm{def:g12:arith:divides}{keterbagian}, \omterm{def:g12:arith:prime}{bilangan prima}, sisa.
Lama dianggap yang paling murni di antara matematika murni, kini ia melindungi
setiap pembayaran daring: sistem sandi RSA bersandar pada teorema Bézout, Gauss
dan Fermat yang dibuktikan di bab ini.

\section{Keterbagian dan pembagian Euklides}

\begin{definition}[Keterbagian]\label{def:g12:arith:divides}
Misalkan $a, b \in \Z$. Kita katakan $b$ \emph{membagi} $a$\index{keterbagian},
ditulis $b \mid a$, jika ada $k \in \Z$ dengan $a = kb$. Kita katakan pula bahwa
$a$ adalah \emph{kelipatan} dari $b$.
\end{definition}

\begin{proposition}\label{prop:g12:arith:divprops}
Jika $c \mid a$ dan $c \mid b$, maka $c$ \omterm{def:g12:arith:divides}{membagi} setiap gabungan bulat
$au + bv$ ($u, v \in \Z$). Jika $a \mid b$ dan $b \mid a$ dengan
$a,b \in \N$, maka $a = b$. Jika $a \mid b$ dan $b \neq 0$, maka $\abs a \leq \abs b$.
\end{proposition}

\begin{proof}
Tulislah $a = kc$, $b = lc$: maka $au + bv = (ku + lv)c$. Butir lainnya
menyusul dari $\abs{a} = \abs{k}\,\abs{b}$ dengan $\abs k \geq 1$ ketika
$b = ka \neq 0$.
\end{proof}

\begin{theorem}[Pembagian Euklides]\label{thm:g12:arith:euclid}
\index{pembagian Euklides}
Misalkan $a \in \Z$ dan $b \in \N^*$. Ada tepat satu pasangan
$(q, r) \in \Z \times \N$ sedemikian sehingga
\[
a = bq + r \qquad\text{dan}\qquad 0 \leq r < b .
\]
Di sini $q$ disebut \emph{hasil bagi} dan $r$ disebut \emph{sisanya}.
\end{theorem}

\begin{proof}
\emph{Keberadaannya.} Himpunan kelipatan $b$ yang tidak melampaui $a$ punya
anggota terbesar $bq$ (himpunan itu tak kosong dan \omterm{def:g12:seq:bounded}{terbatas di atas}); ambil $r = a - bq$.
Menurut kemaksimalannya, $b(q+1) > a$, sehingga $0 \leq r < b$.
\emph{Ketunggalannya.} Jika $bq + r = bq' + r'$ dengan $0 \leq r, r' < b$, maka
$b(q - q') = r' - r$ dan $\abs{r' - r} < b$: kelipatan $b$ yang \omterm{def:g10:numbers:abs}{nilai mutlaknya}
kurang dari $b$ haruslah $0$, sehingga $r = r'$ dan $q = q'$.
\end{proof}

\section{Kekongruenan}

\begin{definition}[Kekongruenan]\label{def:g12:arith:congruence}
Misalkan $n \in \N^*$. Dua \omterm{def:g10:numbers:sets}{bilangan bulat} $a, b$ disebut \emph{kongruen modulo
$n$}\index{kekongruenan}, ditulis $a \equiv b \pmod n$, jika $n \mid (a - b)$
--- setara dengan mengatakan bahwa $a$ dan $b$ bersisa sama pada \omterm{thm:g12:arith:euclid}{pembagian
Euklides} oleh $n$.
\end{definition}

\begin{proposition}[Keselarasan dengan operasinya]\label{prop:g12:arith:congops}
Jika $a \equiv b \pmod n$ dan $c \equiv d \pmod n$, maka
\[
a + c \equiv b + d, \qquad
ac \equiv bd, \qquad
a^k \equiv b^k \ (k \in \N) \pmod n .
\]
\end{proposition}

\begin{proof}
Bilangan $n$ \omterm{def:g12:arith:divides}{membagi} $(a-b) + (c-d) = (a+c) - (b+d)$, dan
$ac - bd = a(c - d) + d(a - b)$ juga kelipatan $n$. Kaidah pemangkatannya
menyusul lewat induksi dari kaidah hasil kalinya.
\end{proof}

\begin{method}[Menghitung pangkat modulo $n$]\label{met:g12:arith:powers}
Untuk menghitung $a^k \bmod n$, susutkan bilangan pokoknya modulo $n$, lalu carilah
pangkat kecil dari $a$ yang kongruen dengan $\pm1$, dan pakailah itu untuk
meruntuhkan eksponennya. Misalnya $2^{100} \bmod 7$: karena $2^3 = 8 \equiv 1 \pmod 7$
dan $100 = 3\times33 + 1$,
\[
2^{100} = \left(2^{3}\right)^{33} \times 2 \equiv 1^{33}\times 2 = 2 \pmod 7 .
\]
\end{method}

\section{FPB, Bézout dan Gauss}

\begin{definition}[Pembagi persekutuan terbesar]\label{def:g12:arith:gcd}
Misalkan $a, b$ \omterm{def:g10:numbers:sets}{bilangan bulat} yang tidak keduanya nol. \emph{Pembagi persekutuan
terbesar}\index{pembagi persekutuan terbesar} $\gcd(a, b)$ adalah \omterm{def:g10:numbers:sets}{bilangan bulat}
terbesar yang \omterm{def:g12:arith:divides}{membagi} $a$ sekaligus $b$. Ketika $\gcd(a,b) = 1$, maka $a$ dan $b$
disebut \emph{saling prima}\index{saling prima}.
\end{definition}

\begin{proposition}[Algoritme Euklides]\label{prop:g12:arith:euclidalgo}
\index{algoritme Euklides}
Jika $a = bq + r$ ($b \neq 0$), maka $\gcd(a, b) = \gcd(b, r)$. Karena itu
melelarkan pembagian Euklidesnya menghitung $\gcd(a,b)$: yaitu sisa tak nol
yang terakhir.
\end{proposition}

\begin{proof}
Setiap pembagi persekutuan $a$ dan $b$ \omterm{def:g12:arith:divides}{membagi} $r = a - bq$
(\cref{prop:g12:arith:divprops}), sehingga menjadi pembagi persekutuan $b$ dan $r$;
begitu pula sebaliknya, sebab $a = bq + r$. Kedua pasangan itu punya pembagi
persekutuan yang sama, jadi punya FPB yang sama. Algoritmenya berhenti sebab sisanya
membentuk \omterm{def:g12:seq:sequence}{barisan} \omterm{def:g10:numbers:sets}{bilangan bulat} tak negatif yang turun tegas.
\end{proof}

\begin{example}\label{ex:g12:arith:euclidalgo}
$\gcd(252, 198)$: $252 = 198 + 54$; $198 = 3\times54 + 36$;
$54 = 36 + 18$; $36 = 2 \times 18 + 0$. Jadi $\gcd(252,198) = 18$.
\end{example}

\begin{theorem}[Identitas Bézout]\label{thm:g12:arith:bezout}
\index{identitas Bézout}
Misalkan $a, b$ \omterm{def:g10:numbers:sets}{bilangan bulat} yang tidak keduanya nol, dan $d = \gcd(a,b)$. Ada
$u, v \in \Z$ sedemikian sehingga
\[
au + bv = d .
\]
Khususnya, $a$ dan $b$ \omterm{def:g12:arith:gcd}{saling prima} jika dan hanya jika $au + bv = 1$ untuk suatu
\omterm{def:g10:numbers:sets}{bilangan bulat} $u, v$.
\end{theorem}

\begin{proof}
Jalankan \omterm{prop:g12:arith:euclidalgo}{algoritme Euklides} secara mundur: setiap sisanya adalah gabungan bulat
dari dua sisa sebelumnya, dan data awalnya $a, b$ adalah gabungan dari
dirinya sendiri; lewat substitusi menurun, sisa tak nol yang terakhir $d$ adalah
gabungan bulat dari $a$ dan $b$. (Pada \cref{ex:g12:arith:euclidalgo}:
$18 = 54 - 36 = 54 - (198 - 3\times54) = 4\times54 - 198 =
4(252 - 198) - 198 = 4\times252 - 5\times198$.)

Untuk kesetaraannya: jika $\gcd(a,b) = 1$, maka Bézout menyediakan $u, v$;
sebaliknya, setiap pembagi persekutuan $a$ dan $b$ \omterm{def:g12:arith:divides}{membagi} $au + bv = 1$, sehingga
memaksa $\gcd(a,b) = 1$.
\end{proof}

\begin{theorem}[Lema Gauss]\label{thm:g12:arith:gauss}
\index{lema Gauss}
Misalkan $a, b, c \in \Z$. Jika $a \mid bc$ dan $\gcd(a, b) = 1$, maka $a \mid c$.
\end{theorem}

\begin{proof}
Bézout memberi $au + bv = 1$; kalikanlah dengan $c$: $acu + bcv = c$. Kedua suku
pada ruas kirinya kelipatan $a$ (yang kedua sebab $a \mid bc$), sehingga
begitu pula $c$.
\end{proof}

\begin{corollary}\label{cor:g12:arith:coprimeprod}
Jika $a \mid c$, $b \mid c$ dan $\gcd(a,b) = 1$, maka $ab \mid c$.
\end{corollary}

\begin{proof}
Tulislah $c = ak$. Dari $b \mid ak$ dan $\gcd(a,b)=1$, Gauss memberi
$b \mid k$, sebutlah $k = bl$; maka $c = abl$.
\end{proof}

\section{Bilangan prima}

\begin{definition}[Prima]\label{def:g12:arith:prime}
\omterm{def:g10:numbers:sets}{Bilangan bulat} $p \geq 2$ disebut \emph{prima}\index{bilangan prima} jika
pembagi positifnya hanyalah $1$ dan $p$.
\end{definition}

\begin{proposition}\label{prop:g12:arith:primedivides}
Setiap \omterm{def:g10:numbers:sets}{bilangan bulat} $n \geq 2$ punya pembagi \omterm{def:g12:arith:prime}{prima}; jika $n$ bukan \omterm{def:g12:arith:prime}{prima}, ia punya
pembagi \omterm{def:g12:arith:prime}{prima} $\leq \sqrt n$. Jika suatu \omterm{def:g12:arith:prime}{prima} $p$ \omterm{def:g12:arith:divides}{membagi} hasil kali $ab$, maka
$p \mid a$ atau $p \mid b$ (\emph{lema Euklides}).
\end{proposition}

\begin{proof}
Pembagi terkecil $d \geq 2$ dari $n$ adalah \omterm{def:g12:arith:prime}{prima} (setiap pembagi sejati $d$
akan menjadi pembagi $n$ yang lebih kecil). Jika $n = de$ tersusun dengan
$2 \leq d \leq e$, maka $d^2 \leq de = n$, sehingga $d \leq \sqrt n$. Untuk lema
Euklidesnya: jika $p \nmid a$, maka $\gcd(p, a) = 1$ (sebab pembagi $p$ hanyalah
$1$ dan $p$), dan lema Gauss memberi $p \mid b$.
\end{proof}

\begin{theorem}[Euklides]\label{thm:g12:arith:infiniteprimes}
Ada tak berhingga banyak \omterm{def:g12:arith:prime}{bilangan prima}.
\end{theorem}

\begin{proof}
Diberikan sebarang daftar berhingga $p_1, \dots, p_k$ berisi \omterm{def:g12:arith:prime}{bilangan prima}, tinjaulah
$N = p_1 p_2 \cdots p_k + 1$. Ada \omterm{def:g12:arith:prime}{prima} $p$ yang \omterm{def:g12:arith:divides}{membagi} $N$; tetapi tak satu pun $p_i$
\omterm{def:g12:arith:divides}{membagi} $N$ (sebab sisanya $1$), sehingga $p$ \omterm{def:g12:arith:prime}{prima} yang tak ada di daftarnya. Jadi tak
ada daftar berhingga yang menghabiskan \omterm{def:g12:arith:prime}{bilangan primanya}.
\end{proof}

\begin{theorem}[Teorema dasar aritmetika]\label{thm:g12:arith:fta}
\index{teorema dasar aritmetika}
Setiap \omterm{def:g10:numbers:sets}{bilangan bulat} $n \geq 2$ adalah hasil kali \omterm{def:g12:arith:prime}{bilangan prima}, dan pemfaktoran
itu tunggal sampai pada urutan faktornya:
\[
n = p_1^{\alpha_1} p_2^{\alpha_2} \cdots p_r^{\alpha_r},
\qquad p_1 < p_2 < \dots < p_r \text{ prima},\ \alpha_i \geq 1 .
\]
\end{theorem}

\begin{proof}
\emph{Keberadaannya}, lewat induksi kuat: $n$ yang \omterm{def:g12:arith:prime}{prima} adalah pemfaktorannya
sendiri; jika tidak, $n = de$ dengan $2 \leq d, e < n$, dan keduanya terfaktorkan
menurut hipotesis induksinya. \emph{Ketunggalannya}: andaikan
$p_1\cdots p_s = q_1 \cdots q_t$ (\omterm{def:g12:arith:prime}{bilangan prima}, pengulangan diperbolehkan). Menurut
lema Euklides, $p_1$ \omterm{def:g12:arith:divides}{membagi} suatu $q_j$, dan karena \omterm{def:g12:arith:prime}{prima}, $p_1 = q_j$;
coretlah lalu ulangi. Kedua pemfaktorannya cocok suku demi suku.
\end{proof}

\begin{theorem}[Teorema kecil Fermat]\label{thm:g12:arith:fermat}
\index{teorema kecil Fermat}
Misalkan $p$ \omterm{def:g12:arith:prime}{prima} dan $a \in \Z$ dengan $p \nmid a$. Maka
\[
a^{p-1} \equiv 1 \pmod p .
\]
Untuk setiap $a \in \Z$ (tanpa anggapan \omterm{def:g12:arith:gcd}{saling prima}), $a^p \equiv a \pmod p$.
\end{theorem}

\begin{proof}
Tinjaulah $p - 1$ \omterm{def:g10:numbers:sets}{bilangan bulat} $a, 2a, 3a, \dots, (p-1)a$ modulo $p$. Tak satu pun
$\equiv 0$ (sebab jika $p \mid ka$ dengan $1 \leq k \leq p-1$, lema Euklides memaksa
$p \mid k$, dan itu mustahil), dan semuanya berbeda sepasang demi sepasang modulo $p$
(sebab jika $ka \equiv la$, maka $p \mid (k - l)a$, sehingga $p \mid k - l$ dan $k = l$).
Karena itu, modulo $p$, semuanya adalah bilangan $1, 2, \dots, p-1$ dalam suatu urutan.
Dengan mengalikan semua \omterm{def:g12:arith:congruence}{kekongruenannya}:
\[
a^{p-1}\,(p-1)! \equiv (p-1)! \pmod p .
\]
Karena $p$ tidak \omterm{def:g12:arith:divides}{membagi} satu pun dari $1, \dots, p-1$, pemakaian berulang lema
Euklides membolehkan pencoretan $(p-1)!$, sehingga tersisa $a^{p-1} \equiv 1$. Bentuk
keduanya menyusul dengan mengalikannya dengan $a$ (dan sepele ketika $p \mid a$).
\end{proof}

\begin{example}[Penerapan pada persandian]\label{ex:g12:arith:rsa}
Teorema Fermat membuat pemangkatan modulo $n$ dapat dibalik ketika
eksponennya dipilih dengan tepat --- dan itulah jantung sistem sandi \emph{RSA}.
Dengan $p, q$ \omterm{def:g12:arith:prime}{bilangan prima} besar dan $n = pq$, orang menerbitkan $n$ dan sebuah
eksponen $e$; penyandiannya $x \mapsto x^e \bmod n$. Pembacaan sandinya menuntut
eksponen $d$ dengan $ed \equiv 1 \pmod{(p-1)(q-1)}$, yang hanya dapat dihitung
oleh orang yang tahu $p$ dan $q$ --- lagi pula memperoleh kembali $p, q$ dari $n$
berarti \omterm{def:g10:algebra:expand}{memfaktorkan} bilangan yang panjangnya ratusan angka, dan tak ada algoritme
yang diketahui melakukannya dalam waktu yang masuk akal.
\end{example}

\section{Latihan}

\begin{exercise}[$\star$]\label{exo:g12:arith:1}
Hitunglah hasil bagi dan sisa \omterm{thm:g12:arith:euclid}{pembagian Euklides} $2026$ oleh
$17$, lalu $-2026$ oleh $17$.
\end{exercise}

\begin{exercise}[$\star$]\label{exo:g12:arith:2}
Berapakah sisa $7^{100}$ modulo $10$? (Berapakah angka terakhir
$7^{100}$?)
\end{exercise}

\begin{exercise}[$\star$]\label{exo:g12:arith:3}
Dengan \omterm{prop:g12:arith:euclidalgo}{algoritme Euklides}, hitunglah $\gcd(1071, 462)$, lalu carilah \omterm{def:g10:numbers:sets}{bilangan bulat}
$u, v$ dengan $1071u + 462v = \gcd(1071, 462)$.
\end{exercise}

\begin{exercise}[$\star$]\label{exo:g12:arith:4}
Tunjukkan bahwa untuk setiap $n \in \Z$, $n^2$ kongruen dengan $0$ atau $1$ modulo
$4$. Simpulkan bahwa \omterm{def:g10:numbers:sets}{bilangan bulat} yang $\equiv 3 \pmod 4$ tak pernah menjadi
jumlah dua kuadrat.
\end{exercise}

\begin{exercise}[$\star\star$]\label{exo:g12:arith:5}
Tunjukkan bahwa untuk semua $n \in \N$, $n(n+1)(2n+1)$ terbagi oleh $6$.
\end{exercise}

\begin{exercise}[$\star\star$]\label{exo:g12:arith:6}
Selesaikan di $\Z$ \omterm{def:g12:arith:congruence}{kekongruenan} $5x \equiv 3 \pmod{11}$.
(Petunjuk: carilah balikan $5$ modulo $11$.)
\end{exercise}

\begin{exercise}[$\star\star$]\label{exo:g12:arith:7}
Selesaikan di $\Z \times \Z$ \omterm{def:g10:algebra:equation}{persamaan} Diofantin
\[
17x - 40y = 1,
\]
lalu lukiskan semua penyelesaian $17x - 40y = 6$.
\end{exercise}

\begin{exercise}[$\star\star$]\label{exo:g12:arith:8}
Tunjukkan bahwa $\sqrt2$ \omterm{ex:g10:numbers:classify}{irasional}, dengan memakai ketunggalan pemfaktoran
\omterm{def:g12:arith:prime}{primanya} (bandingkan eksponen $2$ pada kedua ruas $a^2 = 2b^2$).
\end{exercise}

\begin{exercise}[$\star\star\star$]\label{exo:g12:arith:9}
Misalkan $p$ suatu \omterm{def:g12:arith:prime}{bilangan prima}.
\begin{enumerate}
  \item Tunjukkan bahwa untuk $1 \leq k \leq p - 1$, $p$ \omterm{def:g12:arith:divides}{membagi} $\dbinom{p}{k}$.
        (Petunjuk: pakailah $k\binom pk = p\binom{p-1}{k-1}$, \cref{exo:g12:comb:7},
        dan lema Gauss.)
  \item Simpulkan, lewat induksi pada $a \geq 0$, suatu bukti lain bagi teorema
        kecil Fermat dalam bentuk $a^p \equiv a \pmod p$.
\end{enumerate}
\end{exercise}

\begin{exercise}[$\star\star\star$]\label{exo:g12:arith:10}
\emph{(Soal sisa Tiongkok.)} Carilah semua \omterm{def:g10:numbers:sets}{bilangan bulat} $n$ sedemikian sehingga
\[
n \equiv 2 \pmod 3, \qquad n \equiv 3 \pmod 5, \qquad n \equiv 2 \pmod 7 .
\]
(Petunjuk: selesaikan dua syarat pertamanya, lalu masukkan yang ketiga; koefisien
Bézout akan membantu.)
\end{exercise}

\section{Soal: Sandi rahasia dan angka pemeriksa}

\begin{problem}[{Soal akhir pekan --- kekongruenan menjaga setiap
kode batang dan kartu kredit, dan teorema kecil Fermat menjalankan
kunci bagi rahasia dunia}]
\label{pb:g12:arith:1}
G.\,H.~Hardy membanggakan diri pada tahun 1940 bahwa teori bilangan ``tak ternoda''
oleh penerapan. Delapan puluh tahun kemudian, setiap bunyi pemindai kode batang, setiap
pembayaran kartu kredit dan setiap pesan bersandi membantahnya
--- justru dengan perkakas bab ini: \omterm{def:g12:arith:congruence}{kekongruenan}
(\cref{prop:g12:arith:congops}), balikan Bézout
(\cref{thm:g12:arith:bezout}) dan teorema kecil Fermat
(\cref{exo:g12:arith:9}). Soal ini memeriksa kodenya, membobol
kunci versi mainannya, lalu belajar mengapa kunci yang sungguhan bertahan.

\medskip
\noindent\textbf{Bagian I --- Kelancaran \omterm{def:g12:arith:congruence}{kekongruenan}.}
\begin{enumerate}
  \item Hitunglah $2026 \bmod 7$; lalu angka terakhir
        $7^{100}$ (carilah kitaran pangkat $7$ modulo
        $10$).
  \item Pemangkatan cepat (\cref{met:g12:arith:powers}):
        hitunglah $5^{117} \bmod 13$ (mulailah dari
        $5^2 \equiv -1$).
  \item Selesaikan $3x \equiv 5 \pmod 7$.
  \item Jalankan \omterm{prop:g12:arith:euclidalgo}{algoritme Euklides} pada $(97, 35)$, substitusikan mundur
        untuk mencari \omterm{def:g10:numbers:sets}{bilangan bulat} $u, v$ dengan $97u + 35v = 1$, lalu
        simpulkan balikan $35$ modulo $97$.
  \item Nyatakan dengan tepat kapan $a$ punya balikan modulo $n$, dan
        teorema mana yang menyerahkan balikannya.
\end{enumerate}

\medskip
\noindent\textbf{Bagian II --- Angka pemeriksa.}
\begin{enumerate}[resume]
  \item ISBN-10: kesepuluh angka $d_1 \dots d_{10}$ pada kode
        sebuah buku harus memenuhi
        $10d_1 + 9d_2 + \dots + 2d_9 + 1d_{10} \equiv 0
        \pmod{11}$. Periksalah ISBN yang sungguhan,
        $0\,306\,40615\,2$.
  \item Buktikan bahwa rancangan ISBN itu mendeteksi \emph{setiap}
        galat satu angka: jika satu angkanya berubah sebesar
        $d \not\equiv 0$, maka jumlah berbobotnya berubah sebesar $w d$
        dengan $1 \leq w \leq 10$ --- mengapa itu tak pernah bisa
        $\equiv 0 \pmod{11}$
        (\cref{thm:g12:arith:gauss})?
  \item Buktikan bahwa rancangan itu juga mendeteksi setiap pertukaran dua
        angka bertetangga (yang berbeda). Lalu jelaskan rahasia
        rancangannya: sifat $11$ yang mana yang membuat kedua buktinya jalan,
        dan apa yang bisa keliru dengan \omterm{def:g12:complex:modulus}{modulus} $10$?
  \item Kode batang EAN-13 memberi bobot $1, 3, 1, 3, \dots$ pada angkanya,
        modulo $10$. Hitunglah angka pemeriksa yang melengkapi
        $978\,2940199\,05$. Pertukaran bertetangga mana yang \emph{gagal}
        dideteksi EAN? (Kapan
        $2(a - b) \equiv 0 \pmod{10}$?)
  \item Kartu kredit memakai rancangan Luhn: dari kanan, lipatduakan
        tiap angka kedua (dengan mengurangi $9$ ketika lipatduanya
        melampaui $9$), jumlahkan semuanya, lalu tuntutlah kelipatan
        $10$. Periksalah nomor ujinya,
        $4539\,1488\,0343\,6467$.
  \item Dalam satu kalimat: apa yang dibeli \omterm{def:g12:complex:modulus}{modulus} \omterm{def:g12:arith:prime}{prima} bagi ISBN,
        yang tak dapat dimiliki EAN dan Luhn karena terantai pada $10$?
\end{enumerate}

\medskip
\noindent\textbf{Bagian III --- Kunci Fermat.}
\begin{enumerate}[resume]
  \item Sebuah jebakan sebelum hartanya: hitunglah
        $2^{10} \bmod 341$, simpulkan $2^{340} \bmod 341$ --- lalu
        faktorkanlah $341$. Apa yang dikatakan contoh ini (sebuah
        \emph{\omterm{def:g12:arith:prime}{prima} semu Fermat}) tentang pemakaian teorema
        kecil Fermat sebagai uji keprimaan?
  \item RSA dalam ukuran mini: ambil $p = 3$, $q = 11$, jadi
        $n = 33$ dan $(p-1)(q-1) = 20$; eksponen umumnya
        $e = 3$. Carilah eksponen rahasianya $d$ dengan
        $3d \equiv 1 \pmod{20}$ (cara pertanyaan 4).
  \item Sandikanlah pesan $m = 4$: hitunglah
        $c = m^3 \bmod 33$.
  \item Bacalah sandinya: hitunglah $c^d \bmod 33$ (pakailah
        $c \equiv -2 \pmod{33}$) lalu perolehlah kembali pesannya.
  \item Mengapa pembacaan sandinya selalu berhasil: tunjukkan bahwa
        $m^{21} \equiv m$ baik modulo $3$ maupun modulo $11$
        (teorema kecil Fermat di masing-masing dunianya), lalu simpulkan
        modulo $33$ (\cref{thm:g12:arith:gauss} merekatkan kedua
        \omterm{def:g12:arith:congruence}{kekongruenannya}). Di manakah bentuk khusus $21 = ed$, yaitu
        $1 + 20k$, masuk?
  \item Keamanan kuncinya: semua orang tahu $n$ dan $e$;
        memperoleh kembali $d$ menuntut $(p-1)(q-1)$, sehingga menuntut faktor
        $n$. Bilangan $33$ kita terfaktorkan sekali lihat --- lalu mengapa rancangan
        yang sama, dengan $n$ sepanjang enam ratus angka, melindungi
        bank sedunia? (Satu kalimat tentang ketaksetangkupan antara
        mengalikan dan \omterm{def:g10:algebra:expand}{memfaktorkan}.)
\end{enumerate}

\medskip
\noindent\textbf{Bagian IV --- Yang klasik.}
\begin{enumerate}[resume]
  \item Pencacahan prajurit Tiongkok kuno (bandingkan
        \cref{exo:g12:arith:10}): sejumlah prajurit menyisakan
        $2$ ketika dibariskan dalam kelompok $3$ dan menyisakan $3$
        ketika dibariskan dalam kelompok $5$. Carilah semua cacah yang
        mungkin, lalu jelaskan mengapa jawabannya tunggal modulo $15$.
  \item Bukti satu baris, akhirnya: dari $10 \equiv 1 \pmod 9$,
        buktikan bahwa setiap bilangan kongruen dengan jumlah angkanya
        modulo $9$; dari $10 \equiv -1 \pmod{11}$, turunkan
        kaidah jumlah berselang-seling untuk $11$. (Jilid sekolah
        menengah pertama membuktikan ini dengan aljabar yang gamblang ---
        kagumilah pemampatannya.)
  \item Penutup --- Hardy melawan kode batang: rangkumlah
        perkakas bab ini (aritmetika \omterm{def:g12:arith:congruence}{kekongruenan}, balikan Bézout,
        teorema kecil Fermat, perekatan \omterm{def:g12:complex:modulus}{modulus} yang \omterm{def:g12:arith:gcd}{saling prima}) dan
        di mana masing-masingnya mengunci pada tempatnya dalam soal ini;
        lalu berikan vonis zaman kini atas kata ``tak ternoda'' itu.

\end{enumerate}
\end{problem}
